\section{Conclusion}\label{sec:conclusion}

This thesis replicates and extends Greenwood and Sammon's investigation of the index effect using OMX Stockholm Benchmark changes from July 2007 through June 2026. The complete return sample contains 256 additions and 215 deletions. Addition CAR falls from 8.52\% in the earliest block to 1.50\% in the latest, but the intervening periods do not form a steady decline. Across 2017--2026, the mean remains 2.15\%, with a 95\% interval of [1.01\%, 3.39\%]. The latest-period interval is too wide to establish that the return has become economically negligible.

The Swedish calendar adds an important qualification. Announcements occur much closer to implementation recently, and the first-to-last decline is concentrated after a common first seven daily returns. The initial seven-day means are 1.19\% and 1.80\%, with an imprecise difference. This does not show that shorter lead times caused the decline; it shows why a declining full-event return should not automatically be described as a declining announcement response.

Benchmark-linked holdings change around the events, and trading becomes more concentrated immediately before effectiveness. A recent FI comparison supports the scale of observed fund purchases. Historical coverage and adjustment restrictions nevertheless make the early return-to-demand ratio too fragile for a precise elasticity comparison. Within-firm share classes do not isolate a well-resolved class-specific response, while fixed-cohort returns through ED+20 retain a positive later-decade addition association.

The resulting conclusion is substantive but bounded: positive later Swedish addition returns coexist with concentrated implementation trading, and their apparent evolution depends on timing and measurement. The data establish neither complete disappearance nor a permanent causal price effect of index membership. The economic distinction is between returns accumulated over the full episode, the response during a fixed initial interval, and trading at implementation. In Sweden, these do not evolve as one interchangeable measure.
